\scanpage{44}
\mspage{- 27 -}
\margin{n° 380}

PROPOSITION 3.- \uline{Pour toute fonction} $f$ \uline{holomorphe sur un voisinage du spectre de} $x\in A$, \uline{on a} $\operatorname{Sp}(f(x))=f(\operatorname{Sp}(x))$.

Supposons en effet que $t\in\operatorname{Sp}(x)$. Alors $f(z)-f(t)$ est holomorphe sur un voisinage de $\operatorname{Sp}(x)$ et est nulle au point $z=t$. Il existe donc une fonction $g(z)$ holomorphe sur un voisinage de $\operatorname{Sp}(x)$ telle que $f(z)-f(t)=g(z)(z-t)$. On en déduit que $g(x)(x-te)=f(x)-f(t)e$. (Propositions 1 et 2). Puisque $(x-te)$ n'est pas inversible dans $A$, il en résulte que $f(x)-f(t)e$ n'est pas inversible dans $A$ donc que $f(t)\in\operatorname{Sp}(f(x))$. Réciproquement, si $t\notin f(\operatorname{Sp}(x))$, alors $h(z)=(f(z)-t)^{-1}$ est holomorphe sur un voisinage de $\operatorname{Sp}(x)$. On a donc $h(x)(f(x)-te)=e$, ce qui prouve que $t\notin\operatorname{Sp}_{(f(x))}$.

COROLLAIRE.- \uline{Pour que} $f(x)$ \uline{soit inversible dans} $A$, \uline{il faut et il suffit que} $0\notin f(\operatorname{Sp}(x))$.

PROPOSITION 4.- \uline{Soient} $x$ \uline{un élément de} $A$, $f$ \uline{une fonction holomorphe sur un voisinage} $U$ \uline{de} $\operatorname{Sp}(x)$ \uline{et} $g$ \uline{une fonction holomorphe sur un voisinage} $V$ \uline{de} $f(\operatorname{Sp}(x))$. \uline{Alors} $(g\circ f)(x)=g(f(x))$.

Soient $d\mu$ une mesure à support compact contenu dans $V\cap(\mathbb C-f(\operatorname{Sp}(x))$ et $W$ un voisinage ouvert de $f(\operatorname{Sp}(x))$ tels que pour tout $z\in W$, on ait $g(z)=\frac{d\mu(\theta)}{z-\theta}$. Soient d'autre part $d\rho$ une mesure à support compact contenu dans $U\cap f^{-1}(W)\cap(\mathbb C-\operatorname{Sp}(x))$ et $W^\circ$ un voisinage ouvert de $\operatorname{Sp}(x)$ tels que, pour toute fonction $F$ holomorphe sur $U\cap f^{-1}(W)$ et tout $z\in W^\circ$, on ait
\[
F(z)=\int\frac{F(t)}{z-t}\,d\rho(t).
\]
La fonction $g\circ f$ étant holomorphe sur $U\cap f^{-1}(W)$, on a
\[
(g\circ f)(x)=\int R(x,t)g(f(t))\,d\rho(t).
\]
Mais, pour tout $t$ appartenant au support de $d\rho$, $f(t)\in W$ et par conséquent

\margin{54}
